% ============================================
% Economist Cover Letter – Penn Wharton Budget Model
% ============================================

\documentclass[11pt,letterpaper]{letter}

\usepackage{geometry}
\usepackage{microtype}
\usepackage[hidelinks]{hyperref}

\geometry{left=25mm,right=25mm,top=25mm,bottom=25mm}

\signature{Decory Edwards}
\address{
Department of Economics \\
Johns Hopkins University \\
Baltimore, MD 21218 \\
\texttt{dedwar65@jh.edu}
}

\begin{document}

\begin{letter}{
Hiring Committee \\
Penn Wharton Budget Model \\
University of Pennsylvania \\
Philadelphia, PA
}

\opening{Dear Members of the Hiring Committee,}

I am writing to express my strong interest in the Economist position with the Penn Wharton Budget Model. I am a Ph.D.\ candidate in Economics at Johns Hopkins University, advised by Professors Christopher D.~Carroll, Nicholas W.~Papageorge, and Stelios Fourakis, and I expect to receive my degree in May~2026. My research fields are macroeconomics, wealth and income dynamics, and computational economics, with a focus on quantitative policy analysis in heterogeneous-agent environments.

My research agenda focuses on understanding the determinants of wealth inequality and the mechanisms through which heterogeneity in returns to wealth shapes aggregate outcomes. In my job-market paper, I estimate the degree of heterogeneity in individual portfolio returns using micro-data from the Survey of Consumer Finances and simulate a structural model in which households face idiosyncratic return risk. I show that allowing for persistent heterogeneity in returns substantially improves the model’s ability to match the empirical distribution of wealth, offering a tractable way to connect micro-level return dispersion to macroeconomic inequality.

A central component of this work is a quantitative comparison of revenue-equivalent capital income and wealth taxes. I show that capital income taxation reduces wealth concentration more effectively by disproportionately affecting high-return households, while the two policies generate meaningfully different welfare implications across the return distribution. These results speak directly to the types of policy questions addressed by the Penn Wharton Budget Model, where distributional incidence, fiscal outcomes, and macroeconomic feedback effects are central.

In a second project, I use the Health and Retirement Study to examine the relationship between interpersonal trust and portfolio performance. Building on evidence of a hump-shaped relationship between trust and income, I show that a similar relationship holds for individual returns to wealth, highlighting a behavioral channel through which social attitudes shape saving and investment outcomes. This work also provides new estimates of the persistent component of returns to net worth, which motivate the calibration of heterogeneity in my structural framework.

My training emphasizes computational economics and quantitative policy evaluation. I have extensive experience solving and simulating heterogeneous-agent models, working with large-scale micro data, and producing counterfactual policy analyses. I am particularly drawn to the Penn Wharton Budget Model because of its commitment to transparent, empirically grounded analysis that directly informs public debate on taxation, fiscal policy, and inequality. I would be excited to contribute to the development and application of PWBM tools for evaluating proposed policy reforms.

Thank you very much for your time and consideration. I would welcome the opportunity to discuss my application further.

\closing{Sincerely,}

\end{letter}
\end{document}
